\documentclass[11pt]{article}

\newif\ifBenchTikzExternalize
\BenchTikzExternalizefalse
\newif\ifBenchMintedCache
\BenchMintedCachefalse
\InputIfFileExists{bench-config.tex}{}{}

\newcommand{\benchtexttoken}{ALPHA}
\newcommand{\benchfiguretoken}{4.00}
\newcommand{\benchpreambletoken}{P0}

\usepackage{amsmath}
\usepackage{booktabs}
\usepackage{hyperref}
\ifBenchMintedCache
  \usepackage[cache=true,cachedir=_minted]{minted}
\else
  \usepackage[cache=false]{minted}
\fi

\title{W5 Minted Workload}
\author{FastXeBench}
\date{}

\begin{document}
\maketitle

\section{Purpose}
This workload stresses code highlighting and cache reuse. The text token is \benchtexttoken, and the citation path stays active via \cite{refalpha}.

\section{Code Blocks}
\begin{minted}[fontsize=\small,linenos]{python}
def score(values):
    total = sum(values)
    return total / max(len(values), 1)
\end{minted}

\begin{minted}[fontsize=\small]{json}
{
  "protocol": "v0.2-prime-measure",
  "cache": "candidate",
  "token": "ALPHA"
}
\end{minted}

\begin{minted}[fontsize=\small]{bash}
latexmk -pdfxe main.tex
python bench/scripts/summarize.py --input-root results/raw
\end{minted}

\section{Reference Output}
Equation~\eqref{eq:minted-loss} keeps the auxiliary file path active, while Figure~\ref{fig:minted-bar} gives the figure-edit scenario a concrete target.
\begin{equation}
\label{eq:minted-loss}
L(x) = \sum_{i=1}^{n} \left(x_i - \bar{x}\right)^2 .
\end{equation}

\begin{figure}[h]
\centering
\rule{\benchfiguretoken cm}{0.6pt}
\caption{A tiny figure surrogate whose width changes under the figure-edit scenario.}
\label{fig:minted-bar}
\end{figure}

\begin{table}[h]
\centering
\begin{tabular}{@{}lr@{}}
\toprule
Artifact & Count \\
\midrule
Code blocks & 3 \\
Highlighted lines & 11 \\
\bottomrule
\end{tabular}
\caption{A small table to keep the workload structurally similar to the others.}
\end{table}

\bibliographystyle{plain}
\bibliography{refs}

\end{document}
